\documentclass[
  headings=optiontotocandhead,
  twoside,
  openright,
  numbers=noenddot,
  toc=flat,
  11pt,
  parskip=full,
  listof=flat,
  numbers=noenddot,
  bibliography=totoc,
  a4paper,
  DIV=14,
  headheight=5mm,
  headinclude=true,
]{scrartcl}

% \usepackage{fontspec}
% \usepackage{titlesec}
\usepackage{float}
\usepackage{array}
\usepackage[main=ngerman]{babel} %texlive-langgerman muss installiert sein
\usepackage{graphicx}
\usepackage{subcaption}
\usepackage{xcolor}
\usepackage{relsize}
\usepackage{lastpage}
\usepackage{amsmath,amssymb}

% \input{glyphtounicode}
% \pdfgentounicode=1

\clubpenalty = 1000 % Schusterjungen verhindern
\widowpenalty = 1000 % Hurenkinder verhindern
\displaywidowpenalty = 1000

\graphicspath{{./assets/}{./plotting/}}

\newcommand{\kopfbild}{\voffset7mm\includegraphics[width=25mm]{TU_Logo_RGB}}
\newcommand{\kopftitel}{\sffamily{\textbf{\large{Laborprotokoll Optik}}}}

\usepackage[automark,footsepline,plainfootsepline]{scrlayer-scrpage}
\setkomafont{pageheadfoot}{\normalcolor\footnotesize\scshape}
\setkomafont{pagenumber}{\normalfont\normalsize}
\clearpairofpagestyles
\ihead{\headmark}
\ihead{\kopfbild}
\ohead{\kopftitel}
\ifoot{\smaller{David Koch, Paul Vetter}}
\ofoot{Seite \pagemark/\pageref{LastPage}}

\title{Laborprotokoll Interferometrie}
\author{David Koch, Paul Vetter}
\date{}

\begin{document}
\cleardoubleoddpage

\begin{titlepage} 
\maketitle 
\tableofcontents 
\end{titlepage}

\begin{titlepage} \maketitle \tableofcontents \end{titlepage}
\pagenumbering{arabic}

\section{Aufbau}
Als Lichtquelle wurde ein He-Ne-Laser 155 und als Interferometriesystem ein Projektions-Michelson-Interferometer verwendet. Zur vertikalen und horizontalen Zentrierung im Eingang des Interferometers dienten zwei kippbarejustierbare Spiegel und zwei Blenden. Sowie einige Linsen und ein Objektiv, siehe Tabelle \ref{tab:linsen}.

\begin{table}[H]
  \centering
  \setlength{\tabcolsep}{10pt}
  \begin{tabular}{ |c|c| }
    \multicolumn{1}{c}{Linse}  & \multicolumn{1}{c}{Brennweite (mm)} \\
    M & 0.65\\
    K & 80 \\
    Ob & 160 \\
    Z & 50 \\
    F & 20 \\
  \end{tabular}
  \caption{Verwendete Linsen}
  \label{tab:linsen}
\end{table}
